% !ZP title      = Introduction to Kinematics
% !ZP type       = lecture
% !ZP topic      = Kinematics
% !ZP topicorder = 1
% !ZP order      = 1
% !ZP summary    = Position, distance vs displacement, speed vs velocity, acceleration, and the equations of motion.
% !ZP tags       = motion, suvat, basics, velocity, acceleration

\section*{Introduction to Kinematics}

\textbf{Kinematics} is the part of mechanics that describes \emph{how}
objects move, without yet asking \emph{why}. We build up four ideas in
order: position, displacement, velocity and acceleration.

\subsection*{1. Position}
The \textbf{position} of an object tells us where it is, measured from a
chosen reference point (the \textbf{origin}). In one dimension we write it
as $x$, in metres (m).

\subsection*{2. Distance vs.\ Displacement}
\begin{itemize}
  \item \textbf{Distance} is the total path length travelled. It is a
        scalar and never negative.
  \item \textbf{Displacement} is the straight-line change in position
        \emph{with direction} --- a vector:
\end{itemize}
\begin{equation}
  \Delta x = x_{\text{final}} - x_{\text{initial}}
\end{equation}
\begin{quote}
A runner completing one lap of a 400 m track travels a \textbf{distance}
of 400 m but has a \textbf{displacement} of 0 m.
\end{quote}

\subsection*{3. Speed vs.\ Velocity}
Average velocity uses displacement:
\begin{equation}
  \bar{v} = \frac{\Delta x}{\Delta t}
\end{equation}
The SI unit is metres per second ($\text{m/s}$). Velocity is a vector.

\subsection*{4. Acceleration}
\begin{equation}
  a = \frac{\Delta v}{\Delta t} = \frac{v - u}{t}
\end{equation}
SI unit: $\text{m/s}^{2}$.

\subsection*{5. The Equations of Motion (SUVAT)}
For straight-line motion with \textbf{constant acceleration}:

\begin{center}
\begin{tabular}{ll}
\textbf{Symbol} & \textbf{Meaning} \\
$s$ & displacement (m) \\
$u$ & initial velocity (m/s) \\
$v$ & final velocity (m/s) \\
$a$ & acceleration (m/s$^2$) \\
$t$ & time (s) \\
\end{tabular}
\end{center}

\begin{align}
  v   &= u + at \\
  s   &= ut + \tfrac{1}{2}at^{2} \\
  v^2 &= u^{2} + 2as \\
  s   &= \frac{(u+v)}{2}\,t
\end{align}

Each equation deliberately leaves out one variable --- pick the one that
does not contain the quantity you are missing.

\subsection*{Worked micro-example}
A cyclist starts from rest ($u=0$) and accelerates at
$a = 1.5\ \text{m/s}^2$ for $t = 4\ \text{s}$:
\begin{equation}
  v = u + at = 0 + (1.5)(4) = 6\ \text{m/s}
\end{equation}
